\title{Controlling Text-to-Image Diffusion by Orthogonal Finetuning}

\begin{abstract}
Large-scale text-to-image diffusion models demonstrate remarkable ability to create photorealistic images from textual descriptions. Effectively steering or controlling these powerful models for various downstream applications remains a significant open challenge. To address this, we propose a principled finetuning approach—Orthogonal Finetuning~(OFT)—to adapt text-to-image diffusion models for downstream tasks. Unlike prior techniques, OFT can provably maintain hyperspherical energy, which quantifies pairwise neuron relationships on the unit hypersphere. We observe that this characteristic is vital for retaining the semantic generation capabilities of text-to-image diffusion models. To enhance finetuning stability, we further introduce Constrained Orthogonal Finetuning (COFT), which adds a radius constraint on the hypersphere. We focus on two key finetuning scenarios for text-to-image models: subject-driven generation, where the objective is to produce images of a particular subject from only a few reference images plus a text prompt, and controllable generation, aiming to enable the model to incorporate additional control inputs. Empirically, our OFT framework surpasses existing methods in terms of both image quality and convergence speed.
\end{abstract}

\section{Introduction}

Recent advances in text-to-image diffusion models~\cite{saharia2022photorealistic,ramesh2022hierarchical,rombach2022high} have delivered outstanding performance in text-guided, high-quality image synthesis. However, text prompts alone can still be ambiguous and often insufficient for precise and fine-grained control over generated images. To overcome this limitation, we concentrate on two categories of text-to-image generation tasks in this work:

\begin{itemize}[leftmargin=*,nosep]

    \item \textbf{Subject-driven generation}~\cite{ruiz2023dreambooth}: Starting from just a few images of a subject, the goal is to generate images featuring the same subject in new contexts, guided by a text prompt.
    \item \textbf{Controllable generation}~\cite{zhang2023adding,mou2023t2i}: Given extra control signals (e.g., canny edges, segmentation masks), the objective is to synthesize images conditioned on both the control input and a text prompt.
\end{itemize}

Both tasks fundamentally revolve around the challenge of how to effectively fine-tune text-to-image diffusion models without compromising the generative capabilities acquired during pretraining. We outline the key criteria for an effective fine-tuning approach as: (1) \emph{training efficiency}: involving fewer trainable parameters and reduced numbers of training epochs, and (2) \emph{generalizability preservation}: maintaining high-quality and diverse generative outputs. To achieve this, fine-tuning is commonly performed either by adjusting neuron weights with a small learning rate (\emph{e.g.}, \cite{ruiz2023dreambooth}) or by incorporating a minor module with re-parameterized neuron weights (\emph{e.g.}, \cite{hulora2022,zhang2023adding}). Although straightforward, neither fine-tuning method reliably ensures the retention of pretrained generative performance. Furthermore, there is a lack of systematic insight into developing effective fine-tuning techniques and selecting appropriate hyperparameters such as training epoch counts. A central challenge lies in the absence of a metric to quantify how well the pretrained generative ability is preserved. Current fine-tuning approaches implicitly assume that a smaller Euclidean distance between the fine-tuned model and the pretrained model corresponds to better preservation of pretrained capabilities. For this reason, fine-tuning methods generally adopt very small learning rates. While this assumption might hold in some cases, we contend that the Euclidean distance alone fails to comprehensively capture the extent of semantic preservation; thus, a more structural metric to evaluate differences between the fine-tuned and pretrained models can significantly enhance both the preservation of pretrained performance and the stability of fine-tuning.

Motivated by the empirical finding that hyperspherical similarity effectively encodes semantic information~\cite{liu2017deep,liu2018decoupled,chen2020angular}, we utilize hyperspherical energy~\cite{liu2018learning} to represent the pairwise relational structure among neurons. Hyperspherical energy is defined as the sum of hyperspherical similarities (such as cosine similarity) between all pairs of neurons within the same layer, reflecting the degree of neuron uniformity on the unit hypersphere~\cite{liu2021learning}. We conjecture that a well-finetuned model should exhibit minimal changes in hyperspherical energy relative to the pretrained model. A straightforward approach would be to include a regularization term to maintain constant hyperspherical energy during finetuning; however, this does not guarantee effective minimization of hyperspherical energy differences. To address this, we exploit an invariance property of hyperspherical energy — the pairwise hyperspherical similarity remains unchanged when the same orthogonal transformation is applied to all neurons. Inspired by this invariance, we introduce Orthogonal Finetuning (OFT), which adapts large text-to-image diffusion models to downstream tasks while preserving their hyperspherical energy. The key idea is to learn a layer-wise shared orthogonal transformation of neurons that keeps their pairwise angles intact. OFT can also be interpreted as modifying the canonical coordinate system of neurons within the same layer. Considering that a smaller Euclidean distance between the finetuned and pretrained models indicates better retention of pretrained performance, we additionally propose an OFT variant — Constrained Orthogonal Finetuning (COFT), which restricts the finetuned model to lie within a hypersphere of fixed radius centered on the pretrained neurons.

The reasoning behind why orthogonal transformation is effective for finetuning neurons partially originates from the 2D Fourier transform, where an image can be broken down into magnitude and phase spectra. The phase spectrum, representing the angular relationship between the input and basis, retains the majority of the semantic information. For instance, an image’s phase spectrum combined with a random magnitude spectrum can still reconstruct the original image without losing its semantic content. This observation implies that altering the directions of neurons is fundamental to semantically adjusting the generated image, which is the objective of both subject-driven and controllable generation. Nevertheless, allowing extensive freedom in changing neuron directions tends to degrade the pretrained generative performance. To limit this freedom, we propose preserving the angle between any two neurons, based on the hypothesis that neuron-to-neuron angles play a vital role in encoding the knowledge within neural networks. Building on this insight, it is natural to learn layer-shared orthogonal transformations for neurons in each layer to ensure the hyperspherical energy remains constant.

Additionally, we take inspiration from orthogonal over-parameterized training~\cite{liu2021orthogonal}, which trains classification neural networks from scratch by orthogonally transforming a randomly initialized network. This approach works because a randomly initialized neural network tends to have provably low hyperspherical energy on average, and the objective in \cite{liu2021orthogonal} is to maintain this low hyperspherical energy throughout training (low energy correlates with better classification generalization~\cite{liu2018learning,lin2020regularizing}). The work in \cite{liu2021orthogonal} demonstrates that orthogonal transformations provide enough flexibility to train generalizable neural networks for classification tasks. By contrast, our focus is on finetuning text-to-image diffusion models to enhance controllability and improve downstream generative quality. We highlight two main differences between OFT and \cite{liu2021orthogonal}. First, while \cite{liu2021orthogonal} aims to minimize hyperspherical energy, OFT seeks to maintain the hyperspherical energy level of the pretrained model, thereby preserving the intrinsic structure formed during pretraining. For diffusion model finetuning, reducing hyperspherical energy could disrupt the original semantic arrangements. Second, OFT strives to reduce deviation from the pretrained model, resulting in a constrained version, whereas \cite{liu2021orthogonal} imposes no such constraints. Achieving a balance between flexibility and stability is critical for effective finetuning, and we contend that our OFT framework successfully accomplishes this balance. Our contributions are summarized as follows:

\begin{itemize}[leftmargin=*,nosep]

    \item We introduce a new finetuning technique called Orthogonal Finetuning (OFT) designed to enhance the controllability of text-to-image diffusion models. To boost stability further, we propose a constrained version that restricts the angular deviation from the pretrained model.
    \item Unlike current finetuning approaches, OFT carries out model finetuning while provably maintaining the hyperspherical energy, which we find empirically to be a crucial metric for preserving the generative semantics of the pretrained model.
    \item We apply OFT to two applications: subject-driven generation and controllable generation. Through extensive empirical evaluations, we show that OFT significantly outperforms previous methods in terms of generation quality, convergence speed, and finetuning stability. Additionally, OFT demonstrates superior sample efficiency, converging well with far fewer training images and epochs.
    \item For controllable generation, we propose a novel control consistency metric to assess controllability. The main concept is to infer the control signal from the generated image and then compare it against the original control signal. This metric further confirms the strong controllability achieved by OFT.
\end{itemize}

\section{Related Work}

\textbf{Text-to-image diffusion models}. Significant advancements~\cite{saharia2022photorealistic,ramesh2022hierarchical,rombach2022high,gu2022vector,nichol2022glide} have been achieved in text-to-image generation, largely driven by the rapid progress of diffusion-based generative models~\cite{ho2020denoising,song2021score,song2021denoising,dhariwal2021diffusion} and vision-language representation learning~\cite{su2020vlbert,lu2019vilbert,radford2021learning,wang2021simvlm,li2022blip,li2023blip,alayrac2022flamingo,schuhmann2022laion}. GLIDE~\cite{nichol2022glide} and Imagen~\cite{saharia2022photorealistic} train diffusion models directly in pixel space. GLIDE jointly trains the text encoder with a diffusion prior using paired text-image datasets, whereas Imagen employs a frozen pretrained text encoder. Stable Diffusion~\cite{rombach2022high} and DALL-E2~\cite{ramesh2022hierarchical} train diffusion models within the latent space. Stable Diffusion utilizes VQ-GAN~\cite{esser2021taming} to learn a visual codebook as the latent space, while DALL-E2 uses CLIP~\cite{radford2021learning} to build a joint latent embedding space for both images and text. Besides diffusion models, generative adversarial networks~\cite{reed2016generative,zhang2017stackgan,xu2018attngan,li2019controllable} and autoregressive models~\cite{ramesh2021zero,ding2021cogview,wu2022nuwa,yuscaling} have also been explored for text-to-image generation. OFT is intrinsically a model-agnostic finetuning strategy and can be applied to any text-to-image diffusion model.

\textbf{Subject-driven generation}. To avoid altering the subject, \cite{avrahami2022blended,nichol2022glide} use a user-provided mask as an additional conditioning input. Inversion techniques~\cite{choi2021ilvr,dhariwal2021diffusion,ramesh2022hierarchical,gal2022image} allow modification of the context while keeping the subject unchanged. \cite{hertz2022prompt} supports both local and global editing without requiring input masks. However, these methods do not effectively preserve identity-specific details of the subject. Pivotal Tuning~\cite{roich2022pivotal} finetunes a generator around an initial inverted latent code, adding a regularization term to retain identity. Similarly, \cite{nitzan2022mystyle} learns a personalized generative face prior from a set of images of a single person. \cite{casanova2021instance} can produce various versions of an instance but may lose specific details unique to that instance. DreamBooth~\cite{ruiz2023dreambooth} finetunes the text-to-image diffusion model using a reconstruction loss and a class-specific prior preservation loss, leveraging a customized token and a few subject images. OFT builds on the DreamBooth framework but instead of naive finetuning with a small learning rate, it finetunes the model through orthogonal transformations.

\textbf{Controllable generation}. Image-to-image translation can be considered a type of controllable generation, with earlier approaches primarily utilizing conditional generative adversarial networks~\cite{isola2017image,zhu2017toward,wang2018high,park2019semantic,choi2018stargan}. Diffusion models have also been employed for image-to-image translation tasks~\cite{saharia2022palette,voynov2022sketch,wang2022pretraining}. Recently, ControlNet~\cite{zhang2023adding} introduced a method to steer a pretrained diffusion model by fine-tuning it to incorporate additional control signals, demonstrating impressive results in controllable generation. A similar and concurrent approach, T2I-Adapter~\cite{mou2023t2i}, also fine-tunes a pretrained diffusion model to enhance controllability over generated images. Building upon the same task framework as in \cite{zhang2023adding,mou2023t2i}, we employ OFT to fine-tune pretrained diffusion models, achieving consistently improved controllability while requiring fewer training samples and fewer parameters for fine-tuning. Importantly, OFT adds no extra computational cost during inference.

\textbf{Model finetuning}. Fine-tuning large pretrained models on downstream tasks has grown increasingly prevalent~\cite{devlin2018bert,brown2020language,he2022masked}. Adaptation techniques, such as those in \cite{hulora2022,houlsby2019parameter,pfeiffer2020adapterfusion}, are widely studied within natural language processing as forms of fine-tuning. LoRA~\cite{hulora2022} is closely related to OFT, positing a low-rank structure for additive weight updates during fine-tuning. By contrast, OFT applies a layer-shared orthogonal transformation to update neuron weights multiplicatively, provably maintaining the pairwise angles between neurons within the same layer, which results in enhanced stability.

\section{Orthogonal Finetuning}

\subsection{Why Does Orthogonal Transformation Make Sense?}

We begin by exploring the rationale behind employing orthogonal transformations in fine-tuning text-to-image diffusion models. This inquiry is divided into two parts: (1) why it is beneficial to fine-tune the angle (i.e., direction) of neurons, and (2) why orthogonal transformations are chosen to adjust these angles.

Regarding the first question, our approach is motivated by the empirical findings in \cite{liu2018decoupled,chen2020angular} which indicate that angular feature differences effectively represent the semantic gap. SphereNet~\cite{liu2017deep} demonstrates that training a neural network with all neurons normalized to lie on a unit hypersphere maintains comparable capacity and even enhances generalization, suggesting that neuron directions encapsulate the most crucial information from the data. To further highlight the significance of neuron angles, we conduct a toy experiment depicted in Figure~\ref{fig:toy_ae}, where a standard convolutional autoencoder is trained on a set of flower images. During training, the feature map is generated using the standard inner product ($z$ represents the output element of the convolution kernel $\bm{w}$ and $\bm{x}$ is the input within the sliding window). At test time, we compare three methods to generate the feature map: (a) the inner product used during training, (b) the magnitude of the neurons, and (c) the angular information. The results in Figure~\ref{fig:toy_ae} reveal that angular information of neurons can almost perfectly reconstruct the input images, while neuron magnitudes provide no useful information. It is important to note that the cosine activation (c) is not used during training; training is solely based on the inner product. These findings suggest that neuron angles (directions) primarily encode the semantic content of input images. Consequently, finetuning neuron directions is likely a more effective strategy to modify semantic information in images.

As for the second question, the simplest method to finetune neuron directions is to simultaneously rotate or reflect all neurons within the same layer, which naturally corresponds to an orthogonal transformation. Although using other angular transformations that rotate different neurons by varying angles could provide more flexibility, we find orthogonal transformations strike a favorable balance between flexibility and structural regularity. Additionally, \cite{liu2021orthogonal} shows that orthogonal transformations possess sufficient capacity for learning neural networks. To substantiate this claim, we conduct an experiment illustrating the effective regularization provided by the orthogonality constraint. Using the controllable generation setup from ControlNet~\cite{zhang2023adding}, the results presented in Figure~\ref{fig:ortho} indicate that our standard OFT method performs consistently and achieves precise control after training completes (epoch 20). In contrast, OFT without the orthogonality constraint fails to generate any realistic images and shows no control capability. This experiment confirms the critical role of the orthogonality constraint in OFT.

\subsection{General Framework}

The traditional finetuning approach generally employs gradient descent with a reduced learning rate to update either the entire model or specific layers. Using a small learning rate inherently promotes minimal divergence from the pretrained model. Essentially, standard finetuning aims to train the model while implicitly minimizing $\thickmuskip=2mu \medmuskip=2mu \| \bm{M} - \bm{M}^0\|$, where $\bm{M}$ represents the finetuned model parameters and $\bm{M}^0$ denotes the pretrained parameters. Although this implicit restriction is intuitively reasonable, it may still permit excessive flexibility when finetuning large models. To tackle this, LoRA imposes an additional low-rank constraint on the weight update, specifically $\thickmuskip=2mu \medmuskip=2mu \text{rank}(\bm{M} - \bm{M}^0)=r'$, where $r'$ is chosen to be a small integer. In contrast, OFT introduces a constraint related to pairwise neuron similarity: $\thickmuskip=2mu \medmuskip=2mu \| \text{HE}(\bm{M})-\text{HE}(\bm{M}^0) \|=0$, with $\text{HE}(\cdot)$ denoting the hyperspherical energy of a weight matrix. For illustration, consider a fully connected layer $\thickmuskip=2mu \medmuskip=2mu \bm{W}=\{\bm{w}_1,\cdots,\bm{w}_n\} \in \mathbb{R}^{d \times n}$ where each neuron $\bm{w}_i \in \mathbb{R}^d$ (and $\bm{W}^0$ represents the pretrained weights). The output vector $\thickmuskip=2mu \medmuskip=2mu \bm{z} \in \mathbb{R}^n$ for this layer is computed via $\thickmuskip=2mu \medmuskip=2mu \bm{z} = \bm{W}^\top \bm{x}$, where $\thickmuskip=2mu \medmuskip=2mu \bm{x} \in \mathbb{R}^d$ is the input. OFT can be understood as minimizing the difference in hyperspherical energy between the finetuned and pretrained models:

\begin{equation}\label{eq:oft_principle}
\footnotesize
\min_{\bm{W}} \left\| \text{HE}(\bm{W}) - \text{HE}(\bm{W}^0) \right\| ~~\Leftrightarrow~~ \min_{\bm{W}} \bigg{\|} \sum_{i \neq j} \|\hat{\bm{w}}_i - \hat{\bm{w}}_j\|^{-1} - \sum_{i \neq j} \|\hat{\bm{w}}^0_i - \hat{\bm{w}}^0_j\|^{-1} \bigg{\|}
\end{equation}

Here, $\thickmuskip=2mu \medmuskip=2mu \hat{\bm{w}}_i := \bm{w}_i / \|\bm{w}_i\|$ denotes the normalized $i$-th neuron, and the hyperspherical energy for the layer $\bm{W}$ is defined as $\thickmuskip=2mu \medmuskip=2mu \text{HE}(\bm{W}) := \sum_{i \neq j} \|\hat{\bm{w}}_i - \hat{\bm{w}}_j\|^{-1}$. It is straightforward to see that the minimal attainable value for Eq.~\eqref{eq:oft_principle} is zero. This minimum can be reached whenever $\bm{W}$ differs from $\bm{W}^0$ merely by a rotation or reflection, that is, $\thickmuskip=2mu \medmuskip=2mu \bm{W} = \bm{R} \bm{W}^0$, where $\thickmuskip=2mu \medmuskip=2mu \bm{R} \in \mathbb{R}^{d \times d}$ is an orthogonal matrix (with determinant $+1$ or $-1$ corresponding to rotation or reflection, respectively). This captures the core principle of OFT: finetuning the neural network by learning orthogonal transformations shared across layers to modify neurons. Form

\begin{equation}\label{eq:oft_general}
    \footnotesize
    \bm{z}=\bm{W}^\top\bm{x}=(\bm{R}\cdot\bm{W}^0)^\top\bm{x},~~~\text{subject to}~\bm{R}^\top\bm{R}=\bm{R}\bm{R}^\top=\bm{I}
\end{equation}

Here, $\bm{W}$ represents the OFT-finetuned weight matrix and $\bm{I}$ stands for the identity matrix. The OFT method is depicted in Figure~\ref{fig:block}. In a manner similar to the zero initialization used in LoRA, it is necessary to ensure that OFT fine-tunes the pretrained model exactly starting from $\bm{W}^0$. To accomplish this, we initialize the orthogonal matrix $\bm{R}$ as the identity matrix so that the fine-tuned model begins with the pretrained weights. To maintain the orthogonality of the matrix $\bm{R}$, differential orthogonalization techniques, as described in \cite{liu2021orthogonal,lezcano2019cheap}, can be employed. We will explain how to maintain orthogonality efficiently from a computational perspective.

\subsection{Efficient Orthogonal Parameterization}\label{sect:eff_ortho}

Traditional orthogonalization approaches, such as the Gram-Schmidt process, although differentiable, tend to be computationally intensive in practice~\cite{liu2021orthogonal}. To improve efficiency, we utilize the Cayley parameterization for generating the orthogonal matrix. More precisely, we define the orthogonal matrix as $\thickmuskip=2mu \medmuskip=2mu \bm{R}=(\bm{I}+\bm{Q})(\bm{I}-\bm{Q})^{-1}$, where $\bm{Q}$ is a skew-symmetric matrix satisfying $\thickmuskip=2mu \medmuskip=2mu \bm{Q}=-\bm{Q}^\top$. This method offers computational benefits at a minor cost — the Cayley parameterization only produces orthogonal matrices with determinant $1$, belonging to the special orthogonal group. Fortunately, we observe that this restriction does not impair practical performance. Even with the Cayley transform parameterization of $\bm{R}$, the matrix can remain parameter-inefficient when $d$ is large. To overcome this, we propose modeling $\bm{R}$ as a block-diagonal matrix composed of $r$ blocks, yielding the following structure:

\begin{equation}
    \footnotesize
    \bm{R}=\text{diag}(\bm{R}_1,\bm{R}_2,\cdots,\bm{R}_r)=\begin{bmatrix}
    \bm{R}_1\in \text{O}(\frac{d}{r}) &  &\\
    &\ddots & \\
    & & \bm{R}_r\in \text{O}(\frac{d}{r})
    \end{bmatrix}\in \text{O}(d)
\end{equation}

where $\text{O}(d)$ represents the orthogonal group of dimension $d$, and $\thickmuskip=2mu \medmuskip=2mu \bm{R}\in\mathbb{R}^{d\times d}$ along with $\thickmuskip=2mu \medmuskip=2mu \bm{R}_i\in\mathbb{R}^{d/r\times d/r}, \forall i$ are orthogonal matrices. In the case where $\thickmuskip=2mu \medmuskip=2mu r=1$, the block-diagonal orthogonal matrix reduces to a conventional unconstrained orthogonal matrix. For an orthogonal matrix sized $\thickmuskip=2mu \medmuskip=2mu d\times d$, the parameter count is $\thickmuskip=2mu \medmuskip=2mu d(d-1)/2$, leading to a computational complexity of $\mathcal{O}(d^2)$. For an $r$-block diagonal orthogonal matrix, the parameter count is $\thickmuskip=2mu \medmuskip=2mu d(d/r-1)/2$, which corresponds to a complexity of $\mathcal{O}(d^2/r)$. Optionally, one can share the block matrices to further reduce the parameters, that is, $\thickmuskip=2mu \medmuskip=2mu \bm{R}_i=\bm{R}_j,\forall i \neq j$. This sharing lowers the parameter complexity to $\mathcal{O}(d^2/r^2)$. Despite these strategies to enhance parameter efficiency, it is important to note that the resulting matrix $\bm{R}$ remains orthogonal, ensuring no loss in maintaining hyperspherical energy.

We analyze the parameter efficiency of OFT compared to LoRA. For LoRA with a low-rank parameter $r'$, the count of trainable parameters is $\thickmuskip=2mu \medmuskip=2mu r'(d+n)$. Assuming both $r$ and $r'$ scale with the neuron dimension $d$ (e.g., $\thickmuskip=2mu \medmuskip=2mu r = r' = \alpha d$ with $\thickmuskip=2mu \medmuskip=2mu 0 < \alpha \leq 1$), the parameter complexity of LoRA becomes $\mathcal{O}(d^2 + d n)$, while that of OFT is $\mathcal{O}(d)$. To illustrate the difference concretely, consider a weight matrix of size $\thickmuskip=2mu \medmuskip=2mu 128 \times 128$: LoRA requires $2,048$ trainable parameters for $\thickmuskip=2mu \medmuskip=2mu r' = 8$, whereas OFT only has $960$ trainable parameters with $\thickmuskip=2mu \medmuskip=2mu r = 8$ (without block sharing).

\subsection{Constrained Orthogonal Finetuning}

To further restrict the flexibility of the original OFT, we impose a constraint that the finetuned model stays close to the pretrained model. Specifically, COFT employs the following forward computation:

\begin{equation}\label{eq:coft_general}
    \footnotesize
    \bm{z} = \bm{W}^\top \bm{x} = (\bm{R} \cdot \bm{W}^0)^\top \bm{x}, \quad \text{s.t.} \quad \bm{R}^\top \bm{R} = \bm{R} \bm{R}^\top = \bm{I}, \quad \norm{\bm{R} - \bm{I}} \leq \epsilon
\end{equation}

which enforces both orthogonality and an $\epsilon$-bounded deviation from the identity matrix. The orthogonality can be ensured via the Cayley parameterization described in Section~\ref{sect:eff_ortho}. However, incorporating the $\epsilon$-deviation constraint into the Cayley-parameterized orthogonal matrix is nontrivial. To better understand the Cayley transform, we use the Neumann series approximation for $\thickmuskip=2mu \medmuskip=2mu \bm{R} = (\bm{I} + \bm{Q})(\bm{I} - \bm{Q})^{-1}$, yielding $\thickmuskip=2mu \medmuskip=2

series converges with respect to the operator norm). Consequently, the constraint $\thickmuskip=2mu \medmuskip=2mu\|\bm{R}-\bm{I}\|\leq\epsilon$ can be incorporated within the Cayley transform, leading to the equivalent condition $\thickmuskip=2mu \medmuskip=2mu \|\bm{Q}-\bm{0}\|\leq\epsilon'$, where $\bm{0}$ represents the zero matrix and $\epsilon'$ is a distinct error hyperparameter from $\epsilon$. This revised constraint on $\bm{Q}$ can be straightforwardly enforced using projected gradient descent. To initialize the orthogonal matrix $\bm{R}$ as the identity, we start with $\bm{Q}$ set to the zero matrix. COFT can be interpreted as the fusion of two explicit constraints: minimizing the hyperspherical energy difference and bounding the deviation from the pretrained model. While the latter constraint is often implicitly employed by existing finetuning approaches, COFT explicitly enforces it. Despite OFT’s strong performance, we find that COFT provides superior finetuning stability compared to OFT due to the presence of this explicit deviation constraint. Figure~\ref{fig:eps_coft} illustrates how $\epsilon$ influences COFT’s performance. It is evident that $\epsilon$ regulates the adaptability of finetuning; a larger $\epsilon$ causes the COFT-finetuned model to approximate the OFT-finetuned model, whereas a smaller $\epsilon$ results in behavior more closely aligned with the pretrained text-to-image diffusion model.

\subsection{Re-scaled Orthogonal Finetuning}

We introduce a straightforward extension of OFT by additionally learning a per-neuron magnitude scaling coefficient. This is motivated by the observation that rescaling neurons does not affect the hyperspherical energy, as magnitudes are normalized out. Specifically, the forward pass is defined as $\thickmuskip=2mu \medmuskip=2mu\bm{z}=(\bm{R}\bm{W}^0\bm{D})^\top\bm{x}$\footnote{\textbf{Erratum}: The NeurIPS camera-ready version mistakenly presents the re-scaled OFT forward pass as $\thickmuskip=2mu \medmuskip=2mu\bm{z}=(\bm{D}\bm{R}\bm{W}^0)^\top\bm{x}$. The original implementation is correct, hence results in Appendix~\ref{app:rescaled} remain unaffected.}, where $\thickmuskip=2mu \medmuskip=2mu \bm{D}=\text{diag}(s_1,\cdots,s_n)\in\mathbb{R}^{n\times n}$ is a learnable diagonal matrix with strictly positive diagonal elements $s_1,\cdots,s_n$. Unlike the original OFT forward pass in Eq.~\eqref{eq:oft_general}, where only $\bm{R}$ is updated during training, here both the diagonal matrix $\bm{D}$ and the orthogonal matrix $\bm{R}$ are trainable. This re-scaled OFT enhances OFT’s flexibility with only a negligible increase in parameter count. Our experiments focus on the original OFT to isolate the effectiveness of orthogonal transformation alone, although we note that re-scaled OFT generally performs better (refer to Appendix~\ref{app:rescaled}).

\section{Intriguing Insights and Discussions}

\textbf{OFT is architecture-agnostic}. In principle, OFT can be applied to any neural network architecture. For Transformers, LoRA is commonly applied to attention weights~\cite{hulora2022}. To ensure fair comparison with LoRA, we restrict OFT finetuning to the attention weights in our experiments. Beyond fully connected layers, OFT is also well-suited for finetuning convolutional layers, because the block-diagonal structure of $\bm{R}$ admits meaningful interpretations in convolutional contexts (unlike LoRA). When the number of blocks matches the number of input channels, each block independently transforms a unique neuron channel, analogous to learning depth-wise convolutional kernels~\cite{chollet2017xception}. Conversely, when all blocks in $\bm{R}$ are shared, OFT applies an orthogonal transformation shared across channels. A preliminary investigation of finetuning convolutional layers using OFT is presented in Appendix~\ref{app:convolution}.

\textbf{Relation to LoRA}. LoRA avoids significant shifts in the pretrained weight matrix by adding a low-rank matrix. In comparison, OFT constrains the transformation applied to the pretrained weight matrix to be orthogonal (full-rank), which prevents the transformation from damaging the pretrained information. We can express OFT's forward pass as $\thickmuskip=2mu \medmuskip=2mu \bm{z}=(\bm{R}\bm{W}^0)^\top\bm{x}=(\bm{W}^0+(\bm{R}-\bm{I})\bm{W}^0)^\top\bm{x}$, where $\thickmuskip=2mu \medmuskip=2mu (\bm{R}-\bm{I})\bm{W}^0$ resembles LoRA's low-rank weight modification. Since $\bm{W}^0$ is generally full-rank, OFT also effectively performs a low-rank weight update when $\thickmuskip=2mu \medmuskip=2mu \bm{R}-\bm{I}$ has low rank. Similar to how LoRA includes a rank parameter $r'$, OFT introduces a diagonal block parameter $r$ to reduce the number of parameters that need training. Notably, LoRA and OFT represent two fundamentally different strategies for parameter efficiency. LoRA leverages the low-rank structure to cut down trainable parameters, while OFT approaches this by exploiting sparsity through block-diagonal orthogonality.

\textbf{Reasons for OFT's faster convergence}. As illustrated in Figure~\ref{fig:toy_ae}, the most impactful update for altering semantic meaning is to change neuron directions—precisely the design purpose of OFT. Additionally, OFT can be interpreted as fine-tuning neurons constrained on a smooth hyperspherical manifold, which leads to a more favorable optimization landscape. This observation has also been supported empirically in \cite{liu2021orthogonal}.

\textbf{Why we do not minimize hyperspherical energy}. A major distinction from \cite{liu2021orthogonal} is that our goal is not to minimize hyperspherical energy. In classification tasks, having neurons without redundancy is preferred. Minimizing hyperspherical energy results in neurons being evenly spaced on the hypersphere, which is not a suitable goal for fine-tuning since it might disrupt pretrained information.

\textbf{Balance between adaptability and stability in finetuning}. We identify an inherent trade-off between adaptability and stability. Conventional finetuning offers the greatest adaptability but suffers from poor stability, often leading to model failure. Surprisingly straightforward, OFT achieves a favorable compromise between adaptability and stability by maintaining the pairwise angles between neurons. Additionally, the block-diagonal parameterization can be interpreted as a stronger form of regularization on the orthogonal matrix.

\textbf{No extra inference cost}. In contrast to ControlNet, our OFT framework does not add any inference overhead to the finetuned model. During inference, we simply multiply the learned orthogonal matrix $\bm{R}$ by the pretrained weight matrix $\bm{W}^0$ to obtain an equivalent weight matrix $\thickmuskip=2mu \medmuskip=2mu \bm{W}=\bm{R}\bm{W}^0$. Therefore, the inference speed remains identical to that of the pretrained model.

\section{Experiments and Results}

\textbf{Overall setup}. In our experiments, we employ Stable Diffusion v1.5~\cite{rombach2022high} as the pretrained text-to-image model. To ensure fairness, we randomly select generated images from each method. For subject-driven generation, we largely follow the approach of DreamBooth~\cite{ruiz2023dreambooth}. For controllable generation, we adopt the general protocols of ControlNet~\cite{zhang2023adding} and T2I-Adapter~\cite{mou2023t2i}. To guarantee a fair comparison with LoRA, we apply OFT or COFT exclusively to the same layer where LoRA is employed. Additional results and detailed settings are available in Appendix~\ref{app:settings}.

\subsection{Subject-driven Generation}

\textbf{Experimental setup}. We utilize DreamBooth~\cite{ruiz2023dreambooth} and LoRA~\cite{hulora2022} as baseline methods. All approaches use the identical loss function as in DreamBooth. For both DreamBooth and LoRA, we generally adhere to the original paper’s recommended best hyperparameter settings. Further results are presented in Appendix~\ref{app:settings},\ref{app:coft_oft},\ref{app:more_qualitative},\ref{app:failure}.

\textbf{Stability and convergence during finetuning}. We begin by assessing the finetuning stability and convergence speed for DreamBooth, LoRA, OFT, and COFT. The outcomes are presented in Figure~\ref{fig:teaser} and Figure~\ref{fig:db_convergence}. It is evident that both COFT and OFT can finetune the diffusion model in a stable manner. After 400 iterations, both DreamBooth and OFT variants demonstrate good control, whereas LoRA struggles to maintain the subject identity. At 2000 iterations, DreamBooth begins to produce degraded images, and LoRA fails to generate the yellow shirt correctly (instead producing yellow fur). In contrast, OFT and COFT continue to maintain stable and consistent control over the generated images. These findings confirm the rapid yet stable convergence offered by our OFT framework in subject-driven generation. We emphasize that robustness to the number of finetuning iterations is crucial, as it alleviates the need for carefully tuning iteration counts for different subjects. Both OFT and COFT allow setting a relatively large number of iterations without extensive tuning. For COFT, with an appropriate choice of $\epsilon$, both learning rate and iteration number become straightforward to configure.

\textbf{Quantitative evaluation}. Following the protocol in \cite{ruiz2023dreambooth}, we perform a quantitative evaluation of subject fidelity (using DINO~\cite{caron2021emerging}, CLIP-I~\cite{radford2021learning}), text prompt fidelity (CLIP-T~\cite{radford2021learning}), and sample diversity (LPIPS~\cite{zhang2018unreasonable}). CLIP-I measures the average pairwise cosine similarity between CLIP embeddings of generated and real images. DINO operates similarly to CLIP-I, except that ViT S/16 DINO embeddings are used. CLIP-T computes the average cosine similarity between CLIP embeddings of text prompts and generated images. Additionally, we evaluate the average LPIPS cosine similarity across generated images of the same subject with identical text prompts. Table~\ref{table:dreambooth_final} shows that COFT and OFT significantly outperform DreamBooth and LoRA in the DINO and CLIP-I metrics, while achieving slightly better or comparable results in prompt fidelity and diversity. Each method was finetuned repeatedly on the same pretrained model with 30 different random seeds to reduce variability. The results demonstrate that our OFT framework not only provides improved convergence and stability but also consistently delivers superior final performance.

\textbf{Qualitative comparison}. To gain a clearer insight into the advantages of OFT, we present several randomly selected examples of subject-driven generation in Figure~\ref{fig:dreambooth_final}. For a fair evaluation, all examples are generated from the same finetuned model for each method, ensuring that no text prompts are individually optimized for their final outputs. For each approach, we choose the model that attains the highest validation CLIP scores. From the results shown in Figure~\ref{fig:dreambooth_final}, it is evident that both OFT and COFT exhibit excellent semantic preservation of the subject, whereas LoRA frequently fails to maintain the subject’s identity (e.g., LoRA entirely loses the subject identity in the bowl example). Additionally, OFT and COFT offer significantly more precise control via text prompts, whereas DreamBooth, despite preserving subject identity, often struggles to produce images that align with the text prompt (e.g., the first row of the bowl example). This qualitative comparison highlights that our OFT framework achieves superior control and subject preservation simultaneously. Furthermore, OFT is not sensitive to the number of iterations, allowing it to perform well even with a large number of iterations, unlike DreamBooth and LoRA. Additional qualitative examples are provided in Appendix~\ref{app:more_qualitative}. We also perform a human evaluation in Appendix~\ref{app:human}, which further confirms the advantages of OFT.

\subsection{Controllable Generation}

\textbf{Settings}. We employ ControlNet~\cite{zhang2023adding}, T2I-Adapter~\cite{mou2023t2i}, and LoRA~\cite{hulora2022} as baseline methods. The main paper focuses on three challenging controllable generation tasks: Canny edge to image (C2I) using the COCO dataset~\cite{lin2014microsoft}, segmentation map to image (S2I) on the ADE20K dataset~\cite{zhou2017scene}, and landmark to face (L2F) on the CelebA-HQ dataset~\cite{karras2018progressive,xia2021tedigan}. Each method is used to finetune Stable Diffusion (SD) v1.5 on these datasets for 20 epochs. Additional results are included in Appendix~\ref{app:more_qualitative}, \ref{app:more_control_task}, and \ref{app:failure}.

\textbf{Convergence}. We assess the convergence rate of ControlNet, T2I-Adapter, LoRA, and COFT on the L2F task through both quantitative and qualitative analyses. For the evaluation metric, we calculate the mean $\ell_2$ distance between the control face landmarks and the predicted face landmarks. In Figure~\ref{fig:face_convergence}, we illustrate the face landmark error of the model fine-tuned over various numbers of epochs. It is evident that both COFT and OFT achieve substantially faster convergence. LoRA requires 20 epochs to reach the performance level that our OFT framework attains by the 8th epoch. We highlight that OFT and COFT have a comparable number of trainable parameters to LoRA (significantly fewer than ControlNet), yet they converge much more efficiently than existing approaches. Additionally, the rapid convergence of OFT is further supported by the findings in Figure~\ref{fig:teaser}. The rightmost example in Figure~\ref{fig:teaser} demonstrates that OFT is considerably more data-efficient than ControlNet and LoRA, as it converges effectively using only 5\% of the ADE20K dataset. For qualitative comparisons, we concentrate on OFT, COFT, and ControlNet, given that ControlNet attains the closest landmark error to our methods. The outcomes in Figure~\ref{fig:face_convergence_qual} indicate that both OFT and COFT exhibit stable convergence with the generated face pose progressively aligning with the control landmarks. Unlike our stable and gradual convergence, ControlNet’s controllability abruptly appears after the 8th epoch, mirroring the sudden convergence pattern noted in \cite{zhang2023adding}. This stable convergence behavior makes our OFT framework significantly more user-friendly in practice, as its training dynamics are much smoother and more predictable. Consequently, this stability facilitates the identification of optimal hyperparameters for OFT.

\textbf{Quantitative comparison}. We propose a control consistency metric to assess the effectiveness of controllable generation. The main concept is to derive the control signal from the generated image and then compare it with the original input control signal. For the C2I task, we calculate IoU and F1 score. For the S2I task, we measure mean IoU, mean accuracy, and overall accuracy. For the L2F task, we evaluate the mean $\ell_2$ distance between the control landmarks and the predicted landmarks. Further details on the consistency metrics are provided in Appendix~\ref{app:settings}. For all methods compared, we utilize the best possible hyperparameter configurations. Results shown in Table~\ref{table:control_gen} demonstrate that both OFT and COFT deliver significantly stronger and more precise control than other methods. We note that adapter-based methods (\emph{e.g.}, T2I-Adapter and ControlNet) exhibit slower convergence and yield inferior final outcomes. Compared to ControlNet, LoRA achieves better performance in the S2I task but worse results in the C2I and L2F tasks. Overall, we find that LoRA’s performance ceiling is relatively limited, even with careful hyperparameter tuning. In contrast, the performance of our OFT framework has not yet plateaued, as we observe continued improvement with an increasing number of trainable parameters. We highlight that our quantitative evaluation for controllable generation represents a novel contribution, as it precisely measures the control capabilities of finetuned models (subject to the accuracy of the off-the-shelf segmentation/detection model).

\textbf{Qualitative comparison}. We further conduct a qualitative comparison of OFT and COFT against leading contemporary methods such as ControlNet, T2I-Adapter, and LoRA. The randomly generated images presented in Figure~\ref{fig:control_final} illustrate that OFT and COFT not only produce high-quality, realistic images but also provide precise control. For the S2I task, it is evident that LoRA fails entirely to generate images consistent with the input segmentation map, whereas ControlNet, OFT, and COFT successfully regulate the image generation. Compared to ControlNet, OFT and COFT generate images with higher fidelity, richer details, and more accurate geometric structures while using significantly fewer model parameters. In the C2I task, OFT and COFT can synthesize semantically coherent images based on rough Canny edge inputs, outperforming T2I-Adapter and LoRA, which yield inferior results. For the L2F task, our approach delivers the most precise pose control of generated faces, even under challenging facial poses. Across all three control tasks, we demonstrate that OFT and COFT consistently produce qualitatively superior images relative to state-of-the-art baselines, underscoring the efficacy of our OFT framework for controllable generation. For a more extensive qualitative assessment, additional examples for all three control tasks are provided in Appendix~\ref{app:control_exp}. Furthermore, we illustrate the robustness of OFT on additional control tasks (such as dense pose to human body, sketch to image, and depth to image) in Appendix~\ref{app:more_control_task}.

\section{Concluding Remarks and Open Problems}

Inspired by the insight that angular relationships between neurons play a critical role in visual semantics, we introduce a straightforward yet powerful finetuning approach—orthogonal finetuning (OFT)—for managing text-to-image diffusion models. Our focus is on two specific text-to-image tasks: subject-driven generation and controllable generation. In comparison with current techniques, OFT offers enhanced controllability and greater finetuning stability, while requiring fewer finetuning parameters. Crucially, OFT does not add any extra inference cost, making it suitable for efficient deployment.

OFT also raises several intriguing open questions. Firstly, the orthogonality in OFT is ensured through Cayley parametrization, which necessitates computing a matrix inverse. This slightly constrains OFT's scalability. Although we mitigate this issue with block diagonal parametrization, accelerating this matrix inversion process in a differentiable manner remains a challenge. Secondly, OFT shows promising potential for compositionality, since the orthogonal matrices generated by multiple OFT finetuning procedures can be multiplied to yield another orthogonal matrix. Investigating whether this set of orthogonal matrices retains the knowledge from all downstream tasks is a compelling area for future research. Lastly, the parameter efficiency of OFT heavily depends on the block diagonal structure, which unavoidably introduces certain biases and restricts flexibility. Finding ways to enhance parameter efficiency more effectively and with fewer biases represents a significant open problem. 

\end{document}